
Le \(\lambda\)-anneau \(A\) est caractérisé par la propriété universelle suivante :

\textbf{Proposition 4.2.} \uline{Pour toute} \(K\)-\(\lambda\)-\uline{algèbre} \(B\), \uline{et toute famille} \((b_i)_{i\in I}\) \uline{d'éléments de} \(B\), \uline{il existe un} \(K\)-\(\lambda\)-\uline{homomorphisme et un seul} \(f\) \uline{de} \(A\) \uline{dans} \(B\) \uline{tel que pour tout} \(i\in I\), \(f(X_i)=b_i\).

Pour que \(f\) soit un \(\lambda\)-homomorphisme, il faut que \(f(\lambda^j X_i)=\lambda^j(b_i)\) pour tous \(i,j\)~; ceci définit \(f\) de façon unique. L'homomorphisme ainsi obtenu est bien un \(\lambda\)-homomorphisme, car pour vérifier que \(\lambda_t\circ f=\Ghat(f)\circ\lambda_t\) il suffit de le vérifier sur un système de générateurs de \(A\) sur \(K\), ce qui résulte de la définition de \(f\) et des relations (2.4.2).

4.3. On peut faire la même construction en utilisant les \(\gamma\) au lieu des \(\lambda\). Soit \(A'=K[X_i,\gamma^n X_i]_{i\in I,\ n\ge2}\)~; pour définir la structure de \(\lambda\)-anneau de \(A'\), on se donne un homomorphisme d'anneaux : \(A'\to\Ghat(A')_\nabla\), par :
\begin{equation}\tag{4.3.1}
  \gamma_t(X_i)=1+X_it+\sum_{n\ge2}\gamma^n X_i\mathbin{.}t^n~;
  \qquad
  \gamma_t(\gamma^n X_i)=\gamma'^n(\gamma_t(X_i)).
\end{equation}
On vérifie ensuite la commutativité de (3.7.2), ce qui montre qu'on obtient bien une structure de \(\lambda\)-anneau.

Le \(\lambda\)-anneau \(A'\) vérifie également la propriété universelle 4.2~; il est donc canoniquement isomorphe (en tant que \(\lambda\)-anneau) au \(\lambda\)-anneau \(A\).

\textbf{Définition 4.4.} \uline{Le} \(\lambda\)-\uline{anneau} \(A\) \uline{sera appelé le} \(\lambda\)-anneau libre engendré sur \(K\) par la famille de générateurs \((X_i)_{i\in I}\).

\textbf{Définition 4.5.} \uline{Soient} \(A\) \uline{un} \(\lambda\)-\uline{anneau}, \(J\) \uline{un idéal de} \(A\). \uline{On dit que} \(J\) \uline{est} \uline{un} \(\lambda\)-idéal de \(A\) \uline{si pour tout} \(x\in J\), \uline{et tout} \(n\ge1\), \(\lambda^n(x)\in J\).

Si \(f:A\to B\) est un \(\lambda\)-homomorphisme, le noyau de \(f\) est un \(\lambda\)-idéal.

Soit \(J\) un \(\lambda\)-idéal~; l'homomorphisme composé : \(A\xrightarrow{\ \lambda_t\ }\Ghat(A)\longrightarrow\Ghat(A/J)\) s'annule sur \(J\) et se factorise donc par \(A/J\), qui se trouve ainsi muni 
